% !TeX root = ../main.tex

\chapter{Experimental Methodology}\label{chapter:experimental-methodology}

\section{Research Design}

The evaluation separates three forms of evidence. First, Pixelless VO (PVO) is
compared with a stationary baseline and two pixel-refined reference
configurations. Second, a component ablation holds the stereo frontend,
TEASER++ registration, rigid-flow residual, backend settings, and encoded input
fixed while enabling reciprocal motion-vector agreement and spatial suppression
individually and jointly. Third, offline diagnostics measure track persistence,
stereo support, image-structure coincidence, and the conditions associated with
large local error. The latter measurements explain the operating envelope but
are not inputs to the estimator.

All trajectory tables use full-sequence runs. The reported PVO trajectory uses
the operating point selected by the bounded parameter study described below;
short smoke tests, abandoned variants, and results from non-equivalent
diagnostic builds are excluded. Every reported configuration completed
all eleven sequences; a failed or incomplete run would be reported as missing
rather than replaced by a numeric value.

\section{Dataset}

Experiments use all eleven sequences of the EuRoC MAV dataset, an established
benchmark for visual odometry and SLAM, with its supplied stereo calibration
and ground truth~\cite{burri2016euroc}. The evaluated sequence
aliases, difficulty labels, and frame counts are listed in
\autoref{tab:euroc-evaluation-set}.

Other relevant benchmarks include TUM-VI~\cite{schubert2018vidataset} and the
Monado SLAM Dataset~\cite{deMayoCremersPireMSD}. The choice of EuRoC for the
final evaluation is primarily motivated by two practical requirements of the
present construction: broad usable coverage after stereo rectification and
the availability of depth at both endpoints of temporal correspondences.
The encoder-stereo stage adapted from Passthrough+ operates on rectified
views~\cite{chaurasia2020passthrough} and can recover depth only where both
cameras observe the scene. EuRoC's nearly parallel camera pair is well suited
to the common pinhole projection used here. TUM-VI's wide-angle images and
Monado's varied headset camera arrangements require more adaptation to retain
useful image support in that projection.

A reanalysis of the archived pilot calibrations quantifies the coverage
restriction. Of the calibrated, safe-radius-valid native left-image pixel
centres, \SI{97.0}{\percent} project inside the chosen rectified image on
EuRoC, compared with \SI{51.5}{\percent} on TUM-VI and
\SIrange{58.4}{78.7}{\percent} across the tested Monado camera pairs.
These fractions describe native-to-rectified image coverage. Temporal
candidates outside this domain cannot obtain stereo depth through the
implemented lookup, irrespective of the availability of temporal motion
vectors. Raw image resolution alone does not explain the dataset choice:
the Monado Index trials used $960\times960$ images, exceeding EuRoC's
$752\times480$ resolution.

The preliminary runs also show substantially less usable stereo-temporal
support. After the stereo geometric checks, approximately
\SIrange{74.1}{74.8}{\percent} of temporal candidates had depth at both
endpoints on the two easy Machine Hall pilot sequences, compared with about
\SI{43}{\percent} on TUM-VI room2 and room5 and
\SIrange{6.4}{17.3}{\percent} on the four tested Monado sequences. Such
missing depth prevents candidates from forming the 3D--3D correspondences
required by this frontend. These coverage and support limitations motivated
concentrating the final feasibility evaluation on EuRoC; further adaptation
of rectification and codec-domain stereo to the other camera configurations
lies outside its scope. The pilot observations concern the tested
configurations and do not establish a general limitation of those datasets
or of Passthrough+.

\begin{table}[htbp]
  \centering
  \caption[EuRoC evaluation sequences]{EuRoC sequences used by the full-sequence evaluation.}
  \label{tab:euroc-evaluation-set}
  \begin{tabular}{@{}lllrl@{}}
    \toprule
    Alias & EuRoC sequence & Difficulty & Frames & Camera rate \\
    \midrule
    EMH01 & MH\_01\_easy      & Easy      & 3682 & \SI{20}{\hertz} \\
    EMH02 & MH\_02\_easy      & Easy      & 3040 & \SI{20}{\hertz} \\
    EMH03 & MH\_03\_medium    & Medium    & 2700 & \SI{20}{\hertz} \\
    EMH04 & MH\_04\_difficult & Difficult & 2032 & \SI{20}{\hertz} \\
    EMH05 & MH\_05\_difficult & Difficult & 2273 & \SI{20}{\hertz} \\
    EV101 & V1\_01\_easy      & Easy      & 2912 & \SI{20}{\hertz} \\
    EV102 & V1\_02\_medium    & Medium    & 1710 & \SI{20}{\hertz} \\
    EV103 & V1\_03\_difficult & Difficult & 2149 & \SI{20}{\hertz} \\
    EV201 & V2\_01\_easy      & Easy      & 2280 & \SI{20}{\hertz} \\
    EV202 & V2\_02\_medium    & Medium    & 2348 & \SI{20}{\hertz} \\
    EV203 & V2\_03\_difficult & Difficult & 1921 & \SI{20}{\hertz} \\
    \bottomrule
  \end{tabular}
\end{table}

Encoded stereo and reverse-temporal streams are aligned by timestamp rather than
row number. Where preprocessing produced a one-frame terminal mismatch, the
unpaired terminal timestamp was omitted from the longer stream; no timestamp
shift or hard common-length truncation was applied.

\section{Compared Configurations}

\begin{description}
  \item[Stationary] Baseline that repeats the initial pose.
  \item[TEASER++ only] Motion-vector method with TEASER++ and rigid-flow optimization, without forward-backward agreement or non-maximum suppression.
  \item[Forward-backward] TEASER++-only method with reciprocal agreement and delayed confirmation.
  \item[NMS] TEASER++-only method with age-prioritized suppression of indistinguishable measurements.
  \item[Pixelless VO (PVO)] Proposed codec-only method using reciprocal checks,
    age-prioritized suppression, TEASER++ registration, and rigid-flow
    optimization.
  \item[VOCA MV-first] Frozen reference~\cite{hilscher2026voca}. Codec motion vectors initialize the first
    KLT pass; tracks not recovered in that pass are retried with unseeded optical
    flow.
  \item[VOCA OF-first] Frozen reference. Ordinary KLT optical flow is attempted
    first; lost tracks are retried with motion-vector initialization.
\end{description}

The stationary baseline and the TEASER++-only estimator are distinct
configurations. The VOCA MV-first and OF-first values are taken from their
frozen full-sequence evaluations; diagnostic runs are not substituted for those
trajectories.

\section{Encoding and Preprocessing Configuration}

All streams are encoded with FFmpeg 4.4.2 and \texttt{libx264} in the H.264
High profile with \texttt{yuv420p} output. The enabled partitions are
\texttt{p8x8}, \texttt{p4x4}, and \texttt{i8x8}. For the temporal streams,
the remaining x264 options
are \texttt{keyint=1000}, \texttt{me=umh}, \texttt{merange=64},
\texttt{subme=6}, \texttt{bframes=0}, and \texttt{ref=1}. The encoder uses
two-pass rate control at
\SI{500}{\kilo\bit\per\second}. Images listed by each camera's timestamp CSV
are supplied through FFmpeg's concat demuxer and encoded at 30 frames per
second. This encoded rate controls the container timestamps but does not replace
the EuRoC nanosecond timestamps used by the estimator and evaluation. H.264
fractional motion is preserved using each vector's reported numerator and scale;
no integer rounding is applied.

For every camera, a reverse-temporal stream is made by reversing the same ordered
image list and applying the identical temporal encoder settings. The extracted
reverse vectors are subsequently inverted and remapped to the corresponding
forward-time EuRoC timestamps. Forward and reverse files are required to contain
the same number of frames as the synchronized dataset metadata.
The single-reference configuration limits prediction dependencies within each
encoded stream, but two-pass rate control and reversal of complete sequences
make this preparation offline. No end-to-end real-time claim follows from
\texttt{ref=1} or the absence of B-frames.

Stereo preprocessing first applies the calibration-dependent common projection
described in \autoref{chapter:method}. It then encodes twelve independent
full-frame stereo streams: disparity shifts 0, 8, and 32 pixels crossed with
grid phases $(0,0)$, $(4,0)$, $(0,4)$, and $(4,4)$ pixels. Each source stereo
pair is expanded to an I-P-P triple and encoded at 90 frames per second with
CRF 23. The temporal motion-search and partition settings are retained, while
\texttt{keyint=3}, \texttt{min-keyint=3}, \texttt{scenecut=0},
\texttt{bframes=0}, and \texttt{ref=1} enforce independent triples. The I-P-P
reference structure and shift strategy are adapted from Passthrough+~\cite{chaurasia2020passthrough}; the mosaic layout,
hardware encoder, and subsequent depth processing are not reproduced.

Preprocessing verifies the frame type of every triple, equal timestamp order for
the stereo pair, exact forward/reverse frame counts, the presence of all twelve
shift combinations, and non-empty validity masks. The full command templates are
given in \autoref{appendix:reproducibility}; the immutable evaluation manifest
also records the exact FFmpeg version and hashes of the configurations used.

\section{Estimator Configuration}

All evaluated estimators are visual-only: IMU input and real-time enforcement are
disabled. The proposed method uses a calibrated anchor/stereo pair, an absolute square-root QR
linearization, square-root marginalization, at most seven retained keyframes
plus the current pose after marginalization, and a keyframe connection-ratio threshold of 0.8. With the
minimum-frames-after-keyframe parameter set to 1 and a strict counter test, a
new keyframe can be requested only after more than one frame has elapsed.
The Levenberg--Marquardt iteration-limit parameter is set to 7, with an
initial damping value of $10^{-4}$. Joint-window visual factors use an observation
standard deviation of 4 pixels and a Huber threshold of 0.5 pixels, applied to
the residual norm before normalization. Rigid-flow factors are enabled for
eligible observations in both the anchor and stereo cameras; host and
ineligible observations retain absolute factors. Standard marginalization and
lost-landmark removal
remain active; the additional pre-marginalization filtering and second
optimization pass evaluated in earlier experiments remain disabled.

\autoref{tab:frozen-frontend-parameters} lists the selected PVO frontend thresholds. Before
TEASER++ estimation, at most one stereo-temporal pair per
$40\times40$-pixel cell is retained using the stereo matching cost. TEASER++
uses fixed scale, a \SI{0.1}{\metre} noise bound, GNC-TLS rotation estimation,
and heuristic maximum-clique inlier selection~\cite{yang2020teaser}. PVO enables the
hybrid reciprocal-confirmation policy and indistinguishable-measurement
suppression. At the original operating point, the four component-ablation
configurations differ only in these two switches; their joint use is labelled
``Joint frontend.'' PVO uses the same complete method with selected frontend
thresholds and backend weights.

\begin{table}[htbp]
  \centering
  \caption[Selected PVO frontend parameters]{Selected PVO compressed-domain frontend parameters.}
  \label{tab:frozen-frontend-parameters}
  \begin{tabular}{@{}lr@{}}
    \toprule
    Parameter & Value \\
    \midrule
    Median-filter grid / distance threshold & $16$ px / $5$ px \\
    Stereo reprojection gate & $1.5$ px \\
    Stereo epipolar gate & $0.005$ \\
    Valid stereo depth interval & $(0.333,30]$ m \\
    TEASER++ noise bound & $0.1$ m \\
    Far-depth threshold / reprojection gate & $5$ m / $1.5$ px \\
    Reciprocal cosine / minimum IoU & $0.98$ / $0.15$ \\
    Reciprocal cycle-error gate & $1.5$ px \\
    Single-endpoint tolerance in suppression model & $0.5625$ px \\
    Suppression endpoint distance & $1.125$ px in $L_\infty$ \\
    Confirmation age & 3 observations \\
    \bottomrule
  \end{tabular}
\end{table}

The numerical values above define the selected PVO operating point. A bounded
twelve-candidate study varied reciprocal-agreement gates, suppression distance,
Huber threshold, and observation standard deviation on all eleven sequences.
The selected point gave the strongest typical-sequence trade-off: it had the lowest
median ATE in the original sweep and remained preferable to the local variants
around it. One local variant reduced median RTE-6 slightly but greatly enlarged
the EV103 failure. Since selection and reporting use the same benchmark, the
result is an in-sample operating point rather than a held-out estimate of
generalization. The evaluated executable and configuration are identified by
SHA-256 hashes in \autoref{appendix:reproducibility}.

\section{Metrics and Aggregation}

Absolute trajectory error (ATE) is reported in metres after a single rigid
$\mathrm{SE}(3)$ alignment between the estimated and reference trajectories~\cite{umeyama1991least,sturm2012benchmark}.
Relative translation error (RTE-6) is the translation error in metres over
non-overlapping six-frame intervals~\cite{sturm2012benchmark}. At the nominal camera rate this corresponds
to \SI{0.3}{\second}. Both per-sequence values are root mean square errors
(RMSE), computed over associated positions for ATE and relative-pose intervals
for RTE-6. ATE measures global trajectory consistency, whereas RTE-6
measures short-horizon drift; a trajectory may therefore have small local errors
that accumulate into a larger aligned ATE.

The main tables report every sequence. Across-sequence summaries use the macro
mean and median so that long sequences do not receive additional weight. The
empirical cumulative distribution and the number of sequences completed are
reported alongside these summaries because one failed sequence can dominate an
arithmetic mean.

Qualitative trajectory plots use the same single three-dimensional alignment
as ATE, with scale fixed to one. They show the saved estimator-origin positions
in two orthogonal projections, with equal metric axes and no separate fitting
of the projections. Every saved estimate is retained, including boundary poses
without a valid ground-truth association; available ground truth is shown over
the saved time span. These views reveal path geometry but not the timing or
cause of a failure, which requires the episode analysis below.

\section{Ablation Studies}

The component study contains four configurations at the original operating
point: TEASER++ only, reciprocal forward-backward agreement only,
age-prioritized non-maximum suppression only, and their joint use. Their encoded
inputs, stereo processing, TEASER++ configuration, rigid-flow residual, backend
settings, and evaluated timestamps are identical. ``Joint frontend'' denotes
the complete method at the original parameter settings. ``PVO'' denotes the same
method after selecting reciprocal-agreement thresholds, suppression distance,
observation standard deviation, and Huber threshold. PVO is reported alongside
the four controlled configurations as a parameter-selected variant of the
complete method, rather than an additional component-ablation cell.

A separate paired experiment isolates the backend observation model at the
selected PVO operating point. It keeps the frontend, encoded inputs, stereo
associations, robust registration, window, and loss parameters fixed, and
changes only whether eligible temporal observations use conventional absolute
reprojection factors or the two-camera rigid-flow factors of
\autoref{eq:rigid-flow-residual}. This comparison is reported separately from
the frontend component ablation because it evaluates a different interface in
the estimator.

\section{Track Characteristics}

Track length is the number of frames in which one confirmed landmark ID is observed. Because admission is delayed until a third observation, provisional two-frame candidates that subsequently fail are counted separately rather than entered into the confirmed-track histogram. Results report the full distribution and survival proportions in addition to the mean, since a small long-lived tail can be hidden by a large number of short tracks. Track length is kept separate from trajectory error so that multiview support remains visible even when aggregate pose metrics appear similar.

Spatial support is evaluated offline and is never supplied to the estimator.
Shi-Tomasi corners are detected with quality level $0.01$, minimum distance 7,
and block size 7. An observation is marked as structurally coincident when its
motion-vector support rectangle contains at least one detected corner. The
analysis reports both the fraction of codec regions containing a corner and the
fraction of corners covered by their union. Raw temporal vectors, candidates
with valid stereo depth at both endpoints, reciprocal candidates, TEASER++
survivors, and suppression-retained candidates are reported separately. The
last row also includes newly selected best-direction births, which remain
provisional until reciprocal confirmation; it is therefore not a strict subset
of the preceding reciprocal-consensus row.
Quality levels $0.05$ and $0.10$ test whether the result is retained for
stronger corner responses. These measurements are post-hoc diagnostics of the
selected PVO run and do not affect its trajectory.

Frontend counters are normalized by the evaluated frame count. Stereo support is
reported as the fraction of queries that return a valid match, together with the
availability of depth at consecutive track endpoints. The reported PVO track and
depth distributions are computed from the selected PVO runs. Track statistics from reference
frontends are used only where their trajectory reproduces the corresponding
frozen reference result; non-equivalent diagnostic replays are excluded.

\section{Failure Analysis}

The saved full-sequence trajectories are inspected for large pose updates,
sustained local error, stationary outputs during camera motion, and recovery.
Each selected episode is examined with its preceding and subsequent frames.
Interval boundaries follow the event being examined; they do not divide the
sequences into equal time or distance bins. The selection is retrospective and
illustrative. Contrasting intervals are included to test whether a proposed
explanation also fits successful motion or fails in a different way.

Local comparisons retain the non-overlapping six-frame intervals used by RTE-6.
Their actual timestamps are used, since missing pose samples can change the
duration of a six-frame interval. Adjacent-pose translation and rotation are
also inspected to separate a single discontinuity from sustained error.
Accumulated error after aligning the first associated pose illustrates what
remains after local recovery; it is explicitly distinguished from the
whole-trajectory alignment used for ATE. The archived absolute-residual
ablation and VOCA OF-first reference are compared over identical timestamps
after checking their complete trajectories against the frozen metrics.

Timestamped diagnostics provide raw and selected candidate counts, emitted
frontend IDs, IDs connected to existing landmarks, and optimization iterations.
The connected count precedes creation of new landmarks. A stereo-valid
candidate has passed the preceding filters and has depth at both temporal
endpoints; a zero count at this stage does not isolate which preceding check
failed. Eligibility for a rigid-flow factor likewise does not measure the
number of independent geometric constraints. Full-sequence track lifetimes,
query-success rates, and mean consecutive depth changes remain sequence-level
summaries and are not attributed to individual events without timestamped
evidence.

Raw images, Shi-Tomasi corners, and an inverse-Laplacian sharpness proxy are
used only for offline inspection. The sharpness proxy depends on texture and
illumination as well as blur, so it is not treated as a direct measurement of
motion blur. Stream frame types are checked where an interruption of temporal
MVs is observed. These measurements help locate changes in the input and
processing chain, while the interpretation separates observed events from
possible causes that the recorded data cannot resolve.

\section{Reproducibility}

The PVO implementation is based on commit \texttt{29d09aa0b14e}. The selected
evaluation executable and selected configuration are identified by hashes because
the suppression distance was exposed as an evaluation-time parameter after that
base revision. The evaluation manifest records these hashes, the FFmpeg and Python versions, the
eleven canonical dataset identifiers, and the commands used to regenerate the
tables and plots. Intermediate data are retained outside the thesis source; the
PDF figures included here are immutable exports of that bundle. Missing values
are not imputed, and results from binaries that fail the selected-trajectory
consistency check are omitted.
